\documentclass{article}
\usepackage[utf8]{inputenc}
\usepackage[T1]{fontenc}
\usepackage[french]{babel}
\usepackage{amsmath}
\setcounter{MaxMatrixCols}{20}
\usepackage{amssymb}
\usepackage{booktabs}
\usepackage{array}
\usepackage{geometry}
\geometry{a4paper, margin=2.5cm}
\usepackage{subfiles}
\usepackage{hyperref}
\usepackage{xcolor}

\hypersetup{
    colorlinks=true,
    linkcolor=blue!70!black,
    citecolor=blue!70!black,
    urlcolor=blue!70!black
}

\title{Estimation de pose d'un dirigeable autonome :\\
\large Formulations et modélisations du filtre de Kalman\\
\normalsize Du modèle cinématique minimal au filtrage sur variété (ESKF)}
\author{Grégory Barou}
\date{\today}

\begin{document}

\maketitle
\tableofcontents
\newpage

\subfile{sections/01_introduction.tex}
\subfile{sections/02_conventions.tex}
\subfile{sections/03_situation1_minimal.tex}
\subfile{sections/04_situation2_inertiel.tex}
\subfile{sections/05_situation3_biais.tex}
\subfile{sections/06_invariance_optimisation.tex}
\subfile{sections/07_annexes.tex}

\end{document}